% Elsevier Editorial Manager (EM) build target -- pdflatex-compatible.
%
% EM runs TeX Live 2022 with pdflatex by default and often ignores the
% `%!TEX TS-program = xelatex` directive, so this wrapper deliberately
% avoids XeTeX-only packages (fontspec, unicode-math). Table column
% widths use `calc` + `\linewidth` which are engine-agnostic. The
% companion `main.tex` (in the same directory) uses XeLaTeX for the
% local publication-quality build; both compile the same body.

\documentclass[preprint,11pt,authoryear,a4paper]{elsarticle}

\usepackage[utf8]{inputenc}
\usepackage[T1]{fontenc}
\usepackage{lmodern}       % Latin Modern -- same metrics as Computer Modern

\usepackage{amsmath,amssymb,amsthm}
\usepackage{booktabs}
\usepackage{graphicx}
\usepackage{siunitx}
\usepackage{xurl}
\usepackage{hyperref}
\hypersetup{hidelinks}
\usepackage[capitalise,nameinlink]{cleveref}
\usepackage{xcolor}
\usepackage{longtable}
\usepackage{enumitem}
\usepackage{array}
\usepackage{calc}
\usepackage{multirow}

\usepackage[a4paper,margin=2.5cm]{geometry}
\usepackage{microtype}
\usepackage[font=small,labelfont=bf,labelsep=period]{caption}
\setlength{\parindent}{1.5em}

% Map the handful of glyphs that survive postprocess_tex.py's
% ensuremath step. Use \DeclareUnicodeCharacter (inputenc's native API)
% rather than \newunicodechar so we don't emit "Redefining Unicode
% character" warnings for §, –, †, ‡ that inputenc already handles.
\DeclareUnicodeCharacter{25AB}{\ensuremath{\,\square\,}}  % ▫
\DeclareUnicodeCharacter{03C1}{\ensuremath{\rho}}         % ρ
\DeclareUnicodeCharacter{0394}{\ensuremath{\Delta}}       % Δ
\DeclareUnicodeCharacter{00A7}{\S}                        % §
\DeclareUnicodeCharacter{2013}{--}                        % –
\DeclareUnicodeCharacter{2020}{\textsuperscript{\dag}}    % †
\DeclareUnicodeCharacter{2021}{\textsuperscript{\ddag}}   % ‡

% Pandoc compatibility shims.
\providecommand{\tightlist}{\setlength{\itemsep}{0pt}\setlength{\parskip}{0pt}}
\providecommand{\pandocbounded}[1]{#1}
\providecommand{\textsubscript}[1]{\ensuremath{_{\text{#1}}}}
\providecommand{\real}[1]{#1}

% Editorial Manager only accepts TIFF for the "Figure" item type, but
% pdflatex cannot embed TIFF via \includegraphics.  The typesetter
% inserts the actual figure at production; here we render a placeholder
% box that names the TIFF file, so reviewers know which upload belongs
% at each position.
\newcommand{\emfig}[1]{%
  \fbox{\parbox[c][4.5cm][c]{0.75\linewidth}{%
    \centering
    \textit{Figure to be inserted at production:}\\[6pt]%
    \texttt{#1.tif}%
  }}%
}

\journal{BioSystems}

\begin{document}

\begin{frontmatter}

\title{Robust error-minimization in the genetic code across physicochemical metrics and variant codes: a graph-theoretic analysis in $\mathrm{GF}(2)^6$}

\author[lc]{Paul Clayworth\fnref{eq}}
\author[bs]{Sergey Kornilov\corref{cor}\fnref{eq}}

\affiliation[lc]{organization={Logocentricity Inc.}, addressline={Austin, TX}, country={USA}}
\affiliation[bs]{organization={Biostochastics, LLC}, addressline={Seattle, WA}, country={USA}}

\cortext[cor]{Corresponding author. \texttt{sergey@biostochastics.com}}
\fntext[eq]{Equal contribution.}

\begin{abstract}
The standard genetic code reduces the impact of point mutations, but the robustness of this property across physicochemical metrics, naturally variant codes, and codon-reassignment mechanisms remains incompletely quantified. Embedding the 64 codons in $\mathrm{GF}(2)^6$ represents the hypercube $Q_6$ as a coordinate-dependent subgraph of the encoding-independent single-nucleotide mutation graph $H(3,4)$, and enables continuous $\rho$-interpolation between the two. Under a block-preserving null ($n = 10{,}000$), the standard code is significantly low-cost across four established, code-independent physicochemical distance metrics with partially overlapping content (Grantham $p = 0.006$; Miyata $p < 0.001$; Woese polar requirement $p = 0.003$; Kyte--Doolittle hydropathy $p = 0.001$), and the signal strengthens monotonically as $\rho$ moves $Q_6 \to H(3,4)$. A structure-aware sensitivity analysis under the alignment-derived ProtSub matrix (Jia \& Jernigan 2021) yields the most extreme percentile of any measure tested ($p = 0.0004$; all five p-values pass Bonferroni at $\alpha = 0.05$). Across the 27 NCBI translation tables, near-optimality is preserved: 11 of 12 informative-distance variants retain top-5\% placement after BH--FDR correction. Natural codon reassignments avoid disrupting codon-family connectivity: under $H(3,4)$, observed events are topology-breaking at relative risk $0.32$ versus the candidate landscape (permutation $p \leq 10^{-4}$), robust to alternative topology definitions and base-to-bit encodings. Event-level conditional-logit modelling shows that topology avoidance and local physicochemical cost provide complementary, only weakly correlated signal ($r_s = 0.15$), and that topology adds explanatory value beyond physicochemistry under both $Q_6$ and encoding-independent $H(3,4)$ adjacency. Retrospective reanalysis of nine genome-recoding datasets is consistent with codon-family topology operating as an evolutionary-trajectory constraint distinct from acute engineering fitness. The contribution is the second axis: code evolution is jointly constrained by physicochemical smoothness and codon-family topological integrity, and these two constraints are partly independent.
\end{abstract}

\begin{keyword}
genetic code evolution \sep error minimization \sep $\mathrm{GF}(2)^6$ representation \sep codon reassignment \sep physicochemical distance \sep codon-family topology \sep conditional logit
\end{keyword}

\begin{highlights}
\item Codon-space geometry links genetic-code robustness and reassignment paths
\item Standard and variant codes preserve broad physicochemical error minimization
\item Reassignments are depleted for codon-family topology-breaking moves
\item Conditional-logit models separate topology from physicochemical similarity
\item Synthetic recoding shows boundary conditions for natural-code constraints
\end{highlights}

\end{frontmatter}

\input{manuscript_body_em.tex}

\bibliographystyle{elsarticle-harv}
\bibliography{references}

\end{document}
